\ifajar
  \section{Latihan}
\else
  \section{Bab 1: Overview Mata Kuliah dan Alur Rekayasa}
\fi

Seluruh file yang dipakai pada latihan ini tersedia di
\berkas{sources/chapter01/}, dihitung relatif terhadap akar repositori.
Latihan ini hanya membutuhkan Python 3 dan Git. Docker, Node.js, dan Chrome
hanya diperiksa keberadaannya.

Seluruh skrip hanya menyentuh komputer sendiri. Pemeriksaan port mencoba
memesan tiap port sebentar lalu melepasnya, tanpa mengirim data ke mana pun,
dan percobaan Git berjalan di direktori sementara. Angka detik tiap mahasiswa
akan berbeda, sedangkan pola perbandingannya akan sama.

\ifajar
  \subsection*{Latihan 1: Smoke Test Lingkungan}
\else
  \subsection*{Prosedur 1: Smoke Test Lingkungan}
\fi

Latihan ini memastikan seluruh perkakas keempat belas bab sudah siap, dan
memperlihatkan bahwa kegagalan yang ditemukan sekarang jauh lebih murah
daripada yang ditemukan di tengah praktikum.

\begin{enumerate}
  \item Buat venv bersama bila belum ada, lalu pasang pustaka basis data,
        seperti pada Listing \ref{lst:c01-lat1-jalan}.
  \item Jalankan \berkas{periksa-lingkungan.py}, lalu catat seluruh barisnya.
  \item Perbaiki tiap pemeriksaan yang gagal: pasang perkakas yang kurang,
        nyalakan Docker, atau hentikan proses yang memakai port bab lain.
  \item Jalankan ulang sampai hasilnya hijau, lalu catat apa saja yang
        diperbaiki.
\ifajar
  \item Isi \tabisi{tab:c01-lat1-hasil}. Tabel \ref{tab:c01-lat1-contoh}
        memperlihatkan contoh yang sudah terisi.
\else
  \item Bandingkan hasil pengukuran sendiri dengan Tabel
        \ref{tab:c01-lat1-contoh}.
\fi
  \item Jawab tiga pertanyaan berikut. Pertama, mengapa pemeriksaan Docker
        menanyai server, bukan hanya mencari perintah \texttt{docker}?
        Kedua, bab mana yang tertahan bila port 8080 sudah terpakai, dan
        kapan kegagalannya akan muncul tanpa \textit{smoke test} ini?
        Ketiga, mengapa versi minimum ditulis di skrip, bukan hanya di
        naskah?
\end{enumerate}

\begin{lstlisting}[language=bash, caption={Menyiapkan venv lalu menjalankan \textit{smoke test}, tersedia di \berkas{sources/chapter01/periksa-lingkungan.py}}, label={lst:c01-lat1-jalan}]
cd sources/chapter01
python3 -m venv ../.venv
../.venv/bin/pip install "psycopg[binary,pool]"
# Dua perintah di atas cukup sekali, bila sources/.venv belum ada.
python3 periksa-lingkungan.py

# Keluarannya periksa-lingkungan.py:
#   Pemeriksaan                      Hasil  Keterangan
#   Python 3.11 atau lebih baru      lolos  3.15.2
#   Node.js 20 atau lebih baru       lolos  v22.23.4
#   npm 10 atau lebih baru           lolos  10.9.9
#   Git 2.30 atau lebih baru         lolos  git version 2.43.0
#   Docker 24 atau lebih baru        lolos  server 29.9.3
#   Docker Compose 2 atau lebih baru lolos  5.7.0
#   Chrome atau Chromium             lolos  Google Chrome 157.0.8103.38
#   psycopg di sources/.venv         lolos  3.3.5
#   Port bab 1 sampai 14             gagal  3 terpakai: 5433, 8000, 8080
#     port 5433 milik Bab 3 sedang dipakai proses lain
#     port 8000 milik Bab 4 sedang dipakai proses lain
#     port 8080 milik Bab 11 sedang dipakai proses lain
#
#   Hasil: merah, 1 pemeriksaan gagal
\end{lstlisting}

\begin{table}[!htb]
  \centering
  \caption{Contoh hasil Latihan 1 yang sudah terisi}
  \label{tab:c01-lat1-contoh}
  \begin{tabularx}{\textwidth}{Xll}
    \hline
    \textbf{Pemeriksaan} & \textbf{Hasil} & \textbf{Keterangan} \\
    \hline
    perkakas dan versinya, tujuh baris & lolos & seluruhnya di atas minimum \\
    \hline
    psycopg di venv & lolos & 3.3.5 \\
    \hline
    port bab 1 sampai 14 & gagal & 5433, 8000, 8080 \\
    \hline
    \textbf{hasil} & \textbf{merah} & 1 pemeriksaan gagal \\
    \hline
  \end{tabularx}
\end{table}

\ifajar
\begin{table}[!htb]
  \centering
  \caption{Lembar isian hasil pemeriksaan Latihan 1}
  \label{tab:c01-lat1-hasil}
  \begin{tabularx}{\textwidth}{Xll}
    \hline
    \textbf{Pemeriksaan} & \textbf{Hasil} & \textbf{Keterangan} \\
    \hline
    perkakas dan versinya, tujuh baris & \dots & \dots \\
    \hline
    psycopg di venv & \dots & \dots \\
    \hline
    port bab 1 sampai 14 & \dots & \dots \\
    \hline
    \textbf{hasil} & \dots & \dots \\
    \hline
  \end{tabularx}
\end{table}
\fi

Contoh di atas merah karena ketiga port dipakai proyek lain di laptop yang
sama, bukan karena perkakasnya kurang. Hasil di laptop lain hampir pasti
berbeda. Yang dilaporkan adalah pemeriksaan mana yang gagal dan apa yang
dilakukan sampai hasilnya hijau.

\ifajar
  \subsection*{Latihan 2: Umur Cabang dan Konflik}
\else
  \subsection*{Prosedur 2: Umur Cabang dan Konflik}
\fi

Latihan ini membuktikan bahwa konflik tumbuh lebih cepat daripada umur
cabangnya, sehingga cabang pantas digabung selagi pendek.

\begin{enumerate}
  \item Jalankan \berkas{simulasi-cabang.py}, seperti pada Listing
        \ref{lst:c01-lat2-jalan}, lalu catat keempat barisnya.
  \item Ubah \texttt{BARIS\_PER\_HARI} di \berkas{simulasi-cabang.py} menjadi
        6, jalankan ulang, lalu catat kolom blok konfliknya.
  \item Kembalikan \texttt{BARIS\_PER\_HARI} menjadi 3.
\ifajar
  \item Isi \tabisi{tab:c01-lat2-hasil}. Tabel \ref{tab:c01-lat2-contoh}
        memperlihatkan contoh yang sudah terisi.
\else
  \item Bandingkan hasil pengukuran sendiri dengan Tabel
        \ref{tab:c01-lat2-contoh}.
\fi
  \item Jawab tiga pertanyaan berikut. Pertama, mengapa blok konflik pada 14
        hari jauh lebih banyak daripada baris yang diubah kedua pengembang?
        Kedua, apa yang terjadi pada konflik ketika tiap pengembang bekerja
        dua kali lebih banyak per hari? Ketiga, berapa hari umur cabang
        terpanjang yang masih bebas konflik pada kedua setelan?
\end{enumerate}

\begin{lstlisting}[language=bash, caption={Mengukur konflik menurut umur cabang, tersedia di \berkas{sources/chapter01/simulasi-cabang.py}}, label={lst:c01-lat2-jalan}]
python3 simulasi-cabang.py

# Keluarannya:
#   300 baris, tiap pengembang mengubah 3 baris per hari, git version 2.43.0
#
#   Umur cabang  Diubah A  Diubah B  Diubah keduanya  Blok konflik
#        1 hari         3         3                0             0
#        3 hari         9         9                0             0
#        7 hari        20        19                1             3
#       14 hari        40        39                2            11
#
# Kolom blok konflik dengan BARIS_PER_HARI = 6: 0, 2, 10, 21
\end{lstlisting}

\begin{table}[!htb]
  \centering
  \caption{Contoh hasil Latihan 2 yang sudah terisi}
  \label{tab:c01-lat2-contoh}
  \begin{tabularx}{\textwidth}{Xll}
    \hline
    \textbf{Umur cabang} & \textbf{Konflik, 3 baris per hari} & \textbf{Konflik, 6 baris per hari} \\
    \hline
    1 hari & 0 & 0 \\
    \hline
    3 hari & 0 & 2 \\
    \hline
    7 hari & 3 & 10 \\
    \hline
    14 hari & 11 & 21 \\
    \hline
  \end{tabularx}
\end{table}

\ifajar
\begin{table}[!htb]
  \centering
  \caption{Lembar isian hasil pengukuran Latihan 2}
  \label{tab:c01-lat2-hasil}
  \begin{tabularx}{\textwidth}{Xll}
    \hline
    \textbf{Umur cabang} & \textbf{Konflik, 3 baris per hari} & \textbf{Konflik, 6 baris per hari} \\
    \hline
    1 hari & \dots & \dots \\
    \hline
    3 hari & \dots & \dots \\
    \hline
    7 hari & \dots & \dots \\
    \hline
    14 hari & \dots & \dots \\
    \hline
  \end{tabularx}
\end{table}
\fi

Benih acaknya tetap, sehingga angka latihan ini sama persis pada komputer
mana pun dengan Git versi yang sama. Versi Git yang berbeda dapat menghitung
batas blok konflik sedikit berbeda, sedangkan polanya tetap: konflik tumbuh
jauh lebih cepat daripada umur cabangnya.

\ifajar
  \subsection*{Latihan 3: Kerangka API dan CI yang Hijau}
\else
  \subsection*{Prosedur 3: Kerangka API dan CI yang Hijau}
\fi

Latihan ini menyalakan kerangka API yang dipakai bab-bab berikutnya, lalu
membuktikan bahwa CI menahan perubahan yang melanggar aturan bisnis.

\begin{enumerate}
  \item Jalankan \berkas{ci-lokal.py}, lalu catat ketiga langkahnya beserta
        detiknya.
  \item Nyalakan \berkas{aplikasi.py}, lalu kirim tiga permintaan seperti
        pada Listing \ref{lst:c01-lat3-jalan}. Catat status HTTP tiap
        permintaan.
  \item Ubah \texttt{harga < HARGA\_MINIMUM\_RUPIAH} di \berkas{layanan.py}
        menjadi \texttt{harga < 0}, jalankan \berkas{ci-lokal.py}, lalu catat
        langkah tempatnya berhenti dan nama uji yang gagal.
  \item Kembalikan \berkas{layanan.py}, lalu pastikan CI kembali hijau.
  \item Buat repositori tim, salin direktori ini ke dalamnya, salin
        \berkas{ci.yml} ke \berkas{.github/workflows/}, sesuaikan
        \texttt{paths}-nya, lalu commit.
\ifajar
  \item Isi \tabisi{tab:c01-lat3-hasil}. Tabel \ref{tab:c01-lat3-contoh}
        memperlihatkan contoh yang sudah terisi.
\else
  \item Bandingkan hasil pengukuran sendiri dengan Tabel
        \ref{tab:c01-lat3-contoh}.
\fi
  \item Jawab tiga pertanyaan berikut. Pertama, mengapa uji API tidak
        dijalankan pada jalan yang merah, padahal salah satu ujinya juga
        akan gagal? Kedua, lapisan mana yang berubah bila SQLite diganti
        PostgreSQL? Ketiga, mengapa satu skrip CI yang sama dipakai di
        laptop dan di server CI?
\end{enumerate}

\begin{lstlisting}[language=bash, caption={CI lokal lalu tiga permintaan ke API, tersedia di \berkas{sources/chapter01/}}, label={lst:c01-lat3-jalan}]
python3 ci-lokal.py

# Keluarannya:
#   Langkah              Hasil   Detik  Keterangan
#   Periksa sintaks      lolos    0.04  seluruh file
#   Uji unit service     lolos    0.06  10 uji
#   Uji API lewat HTTP   lolos    0.10  6 uji
#
#   Hasil: hijau, seluruh langkah lolos

python3 aplikasi.py &
curl -s -X POST http://127.0.0.1:8050/buku \
  -H 'Content-Type: application/json' \
  -d '{"judul": "Basis Data Terapan", "harga": 85000, "stok": 4}'
curl -s -X POST http://127.0.0.1:8050/buku \
  -H 'Content-Type: application/json' \
  -d '{"judul": "Struktur Data", "harga": 0, "stok": 2}'
curl -s http://127.0.0.1:8050/buku/9

# Keluarannya, dengan status HTTP di ujung tiap baris:
# {"nomor": 1, "judul": "Basis Data Terapan", "harga": 85000, "stok": 4} 201
# {"pesan": "harga harus bilangan bulat positif"} 400
# {"pesan": "buku nomor 9 tidak ada"} 404
\end{lstlisting}

\begin{table}[!htb]
  \centering
  \caption{Contoh hasil Latihan 3 yang sudah terisi}
  \label{tab:c01-lat3-contoh}
  \begin{tabularx}{\textwidth}{Xll}
    \hline
    \textbf{Pengamatan} & \textbf{Kode utuh} & \textbf{Kode dirusak} \\
    \hline
    hasil CI & hijau, 0,20 detik & merah, berhenti di uji unit \\
    \hline
    uji yang gagal & tidak ada & \texttt{test\_harga\_nol\_ditolak} \\
    \hline
    status buku sah, harga nol, nomor 9 & 201, 400, 404 & 201, 201, 404 \\
    \hline
  \end{tabularx}
\end{table}

\ifajar
\begin{table}[!htb]
  \centering
  \caption{Lembar isian hasil pengamatan Latihan 3}
  \label{tab:c01-lat3-hasil}
  \begin{tabularx}{\textwidth}{Xll}
    \hline
    \textbf{Pengamatan} & \textbf{Kode utuh} & \textbf{Kode dirusak} \\
    \hline
    hasil CI & \dots & \dots \\
    \hline
    uji yang gagal & \dots & \dots \\
    \hline
    status buku sah, harga nol, nomor 9 & \dots & \dots \\
    \hline
  \end{tabularx}
\end{table}
\fi

Kolom kode dirusak pada baris terakhir menunjukkan akibat yang ditahan CI.
Pada kode yang dirusak, permintaan buku seharga nol dijawab 201, artinya buku
gratis benar-benar masuk katalog. Karena itu \textit{route} dan
\textit{repository} tidak cukup diuji sendiri-sendiri: aturan yang salah di
\textit{service} tetap menghasilkan HTTP yang tampak sehat.
